% !TEX root = ../../main.tex
%==================================================================
% chapters/theory/theorem_h_off_koszul_platonic.tex
%
% Theorem H boundary strata: Koszul locus, positive collision defects,
% singular quotients, and critical level.
%==================================================================

\chapter{Theorem H: family support, singular quotients, and critical level}
\label{chap:theorem-h-off-koszul-platonic}

Theorem~H concerns the curve-level chiral Hochschild complex
\[
 RHH_{\mathrm{ch}}(\cA)
 \;=\;
 R\mathrm{Hom}_{\cA^e}(\cA,\cA)
\]
of a chiral algebra on a fixed smooth projective reference curve~$X$.
Topological Hochschild homology, categorical Hochschild cochains,
formal-disk mode cochains, and product-ambient Hall cochains are four
companion functors with their own source categories.  The comparison
convention is
Convention~\ref{conv:three-hochschild}; the local construction is
Remark~\ref{rem:three-hochschild-discipline}.
The coHochschild statement transports this chiral Hochschild complex
across the complete bimodule--bicomodule equivalence of
Theorem~\ref{thm:ordered-HH-coHH-cohomology}, with
$C=\barB^{\mathrm{ch}}(\cA)$ and diagonal bicomodule~$C_\Delta$.

A family datum $H_H(\cA;S)$ consists of a complete chart model
$Q_\cA$, a quasi-isomorphism
$\gamma_\cA\colon Q_\cA\xrightarrow{\sim}
RHH_{\mathrm{ch}}(\cA)$, and a strong deformation retract onto a
complex $K_{\cA,S}$ supported in a finite set
$S\subset\mathbb Z$.  Theorem~H gives
\[
 \ChirHoch^n(\cA)\cong H^n(K_{\cA,S}),
 \qquad
 \operatorname{Supp}\ChirHoch^\bullet(\cA)\subseteq S.
\]
PBW Koszulness supplies an associated-graded quadratic contraction;
the family datum additionally supplies the OPE-residue lift, incidence
compatibility, completion, and ordered-to-symmetric comparison.
Bakalov--De Sole--Kac compute bounded support $\{0,2,3\}$ for
$\operatorname{Vir}_c$ and dimensions $(2,1,0,\ldots)$ for the
rank-one even superboson.  A bounded-to-chart quasi-isomorphism
transports either result to $RHH_{\mathrm{ch}}$.  De Sole--Kac
variational Poisson cohomology is nonvanishing in unboundedly many
degrees for Virasoro-type Poisson vertex algebras \cite{BDSK21}, so
the concentration statement lives strictly on the stated
Koszul/bounded locus, with the bounded-to-chart comparison part of
the hypothesis package.

\section{Object firewall}
\label{sec:theorem-h-object-firewall}

The objects entering Theorem~H have different types:
\[
\begin{array}{c|c}
\text{object} & \text{role}\\
\hline
\cA & \text{input chiral algebra on }X\\
\barB^{\mathrm{ch}}(\cA) & \text{bar coalgebra}\\
\Omega^{\mathrm{ch}}\barB^{\mathrm{ch}}(\cA) &
 \text{bar--cobar resolution of }\cA\\
\cA^{i}=C_X(s^{-1}V,s^{-2}R) &
 \text{quadratic coalgebra of the chosen presentation}\\
\cA^!=\mathbb D_X(\cA^i) &
 \text{Verdier-dual Koszul algebra}\\
RHH_{\mathrm{ch}}(\cA) &
 \text{chiral Hochschild cochain complex}\\
Z^{\mathrm{der}}_{\mathrm{ch}}(\cA)=RHH_{\mathrm{ch}}(\cA) &
 \text{derived chiral centre, chain-level closed-sector object}\\
\mathrm{Coext}_{C^e}^{\bullet}(C_\Delta,C_\Delta) &
 \text{coHochschild self-extension, after }C=\barB^{\mathrm{ch}}(\cA)
\end{array}
\]
The bar--cobar counit
\[
 \Omega^{\mathrm{ch}}\barB^{\mathrm{ch}}(\cA)\longrightarrow \cA
\]
is the inversion map.  A chosen quadratic presentation gives
$\cA^i=C_X(s^{-1}V,s^{-2}R)$ and the comparison
$q_\cA\colon\cA^i\to\barB^{\mathrm{ch}}(\cA)$; Verdier duality then
forms $\cA^!=\mathbb D_X(\cA^i)$.  Theorem~H begins instead with the
diagonal derived endomorphism complex and a family retract.
All quotients and subcomplexes below are taken in the conformal-weight
completed holonomic $\mathcal D_X$-module derived category. Thus
$F^1_{\mathrm{FM}}$ denotes the closed positive-depth term in the
completed collision-depth filtration, and the depth-zero quotient is a
completed quotient.  Hausdorff completion fixes the ambient topology
for infinite standard families.
The equality
$Z^{\mathrm{der}}_{\mathrm{ch}}(\cA)=RHH_{\mathrm{ch}}(\cA)$ names the
chiral derived centre of the boundary chiral algebra.  A physical
bulk, a categorical centre, or a Hall centre enters through a named
comparison theorem with its source and target categories fixed.

\begin{definition}[Positive collision-depth complex;
\ClaimStatusDefinitional]
\label{def:koszul-defect-complex}
\index{Koszul defect complex}
\index{KDH@$\mathcal K_{\mathrm{DH}}^\bullet$|see {Koszul defect complex}}
\index{chiral Hochschild!Koszul defect}
Let $F^\bullet_{\mathrm{FM}}RHH_{\mathrm{ch}}(\cA)$ be the complete
collision-depth filtration of the bigraded chiral Hochschild complex
used in Lemma~\ref{lem:hochschild-shift-computation} and
Theorem~\ref{thm:hochschild-concentration-E1}. The
\emph{positive collision-depth complex} is the closed subcomplex
\[
 \mathrm{KD}_{\mathrm H}^{\bullet}(\cA)
 \;:=\;
 F^1_{\mathrm{FM}}RHH_{\mathrm{ch}}(\cA).
\]
We also write
\[
 \mathcal K_{\mathrm{DH}}^{\bullet}(\cA)
 := \mathrm{KD}_{\mathrm H}^{\bullet}(\cA)
\]
for the same object in the theorem-spine concordance.  It is the part
of the chiral Hochschild complex supported on collision strata of
positive depth.  Its relation to a proposed support set is determined
by the long exact sequence of
Theorem~\ref{thm:theorem-h-off-koszul-explicit-correction}; acyclicity
of this filtration term is one possible input to a family retract.
\end{definition}

\begin{definition}[Positive-depth Hilbert function;
\ClaimStatusDefinitional]
\label{def:koszul-defect-measure}
\index{Koszul defect measure}
For finite conformal-weight pieces set
\[
 \delta_{\mathrm H}^{(n)}(\cA)(q)
 \;:=\;
 \dim_q H^{n}\!\bigl(\mathrm{KD}_{\mathrm H}^{\bullet}(\cA)\bigr).
\]
The series
\[
 \Delta_{\mathrm H}(\cA)(t,q)
 \;:=\;
 \sum_{n\in\mathbb Z}\delta_{\mathrm H}^{(n)}(\cA)(q)t^n
\]
is defined when the displayed graded dimensions are finite and only
finitely many negative cohomological degrees occur.  Given an upper
bound~$m$ for the depth-zero quotient, its terms in degrees
$n\geq m+1$ govern the upper tail of $RHH_{\mathrm{ch}}(\cA)$.
\end{definition}

\begin{remark}[Role of the positive-depth complex]
\label{rem:koszul-defect-measure-right}
The family proof lifts the Shelton--Yuzvinsky--Priddy contraction of
the quadratic Arnold factor to the residue-twisted coefficient
complex and then totalizes the lift over collision incidences.  The
complex $\mathrm{KD}_{\mathrm H}^{\bullet}(\cA)$ is the positive-depth
term on which this lift acts.  The support set is fixed after the
depth-zero quotient and the connecting maps are computed together.
\end{remark}

\begin{definition}[Collision obstruction pages;
\ClaimStatusDefinitional]
\label{def:theorem-h-collision-obstruction-pages}
\index{Theorem H!collision obstruction}
Fix an integer~$m$ and let
\[
 E_1^{r,s}(\cA)
 =
 H^{r+s}\!\left(
  \operatorname{gr}^{\,r}_{F_{\mathrm{FM}}}
  RHH_{\mathrm{ch}}(\cA)
 \right)
\]
be the collision-depth spectral sequence of
$RHH_{\mathrm{ch}}(\cA)$, with $r$ the collision depth. The first
collision page above the depth-zero cutoff is
\[
 \mathfrak{o}_{\mathrm H,1}^{>m}(\cA)
 :=
 \bigoplus_{\substack{r\geq1\\ r+s>m}}
 E_2^{r,s}(\cA),
\]
where $d_1$ is the residue differential along the first collision
divisors. The permanent positive-depth page is
\[
 \mathfrak{p}_{\mathrm H,\infty}^{>m}(\cA)
 :=
 \bigoplus_{\substack{r\geq1\\ r+s>m}}
 E_\infty^{r,s}(\cA)
 \subset
 \operatorname{gr}_{F_{\mathrm{FM}}}
 H^{>m}\!\bigl(RHH_{\mathrm{ch}}(\cA)\bigr).
\]
Higher differentials determine which $E_2$ classes survive.  When the
depth-zero quotient lies in $D^{\leq m}$, the exact terminal tail is
\[
 \mathfrak{o}_{\mathrm H,\infty}^{>m}(\cA)
 :=
 \mathfrak{o}_{\mathrm H}^{m+1}(\cA)
 \oplus
 \bigoplus_{n\geq m+2}
 H^n\!\bigl(\mathrm{KD}_{\mathrm H}^{\bullet}(\cA)\bigr),
\]
where the first cokernel
$\mathfrak{o}_{\mathrm H}^{m+1}(\cA)$ is defined by the boundary map in
Theorem~\ref{thm:theorem-h-off-koszul-explicit-correction}. Thus the
page-level object $\mathfrak{p}_{\mathrm H,\infty}^{>m}$ determines
the exact tail after the first connecting quotient is formed.
For the specialization $m=2$ we retain the abbreviations
\[
 \mathfrak{o}_{\mathrm H,1}^{\geq3}
 :=\mathfrak{o}_{\mathrm H,1}^{>2},
 \qquad
 \mathfrak{p}_{\mathrm H,\infty}^{\geq3}
 :=\mathfrak{p}_{\mathrm H,\infty}^{>2},
 \qquad
 \mathfrak{o}_{\mathrm H,\infty}^{\geq3}
 :=\mathfrak{o}_{\mathrm H,\infty}^{>2}.
\]
\end{definition}

\section{Family support on a PBW Koszul chart}
\label{sec:theorem-h-on-koszul}

\begin{theorem}[Theorem~H on a PBW Koszul chart;
\ClaimStatusConditional]
\label{thm:theorem-h-on-koszul-locus}
\index{Theorem H!on Koszul locus}
\textup{[Type signature: Open quadrant; completed curve-level chain
presentation; Beilinson level~$3$; hypothesis package
$H_H(\cA;S)$.]}

Let $\cA$ be a PBW chiral Koszul algebra on a smooth projective curve
$X$ and let $S\subset\mathbb Z$ be finite.  Assume the family datum
$H_H(\cA;S)$ of
Definition~\ref{def:theorem-h-collision-depth-package}.  Then
\[
 \ChirHoch^n(\cA)\cong H^n(K_{\cA,S}),
 \qquad
 \operatorname{Supp}\ChirHoch^\bullet(\cA)\subseteq S.
\]
If the cohomology of $K_{\cA,S}$ has finite conformal-weight pieces,
then
\begin{equation}
\label{eq:theorem-h-on-koszul}
 P_{\cA}(t,q)
 \;=\;
 \sum_{n\in S}\dim_q H^n(K_{\cA,S})t^n.
\end{equation}
For a translation-equivariant vertex algebra~$V$, a
quasi-isomorphism $\chi_V^{\mathrm{bd}}$ transports the
Bakalov--De Sole--Kac bounded calculation to the chart.  It gives
$S=\{0,2,3\}$ with one-dimensional groups for
$V=\operatorname{Vir}_c$, and $S=\{0,1\}$ with dimensions $(2,1)$
for a rank-one even superboson.
\end{theorem}

\begin{proof}
The homotopy identities in $H_H(\cA;S)$ identify $Q_\cA$ with
$K_{\cA,S}$.  Composition with $\gamma_\cA$ gives the first
display, and taking graded dimensions gives
\eqref{eq:theorem-h-on-koszul}.  The two benchmarks are
\cite[Theorems~7.2 and~7.4]{BDSK21}; the assumed map
$\chi_V^{\mathrm{bd}}$ transports their cohomology to the chart.
\end{proof}

\begin{remark}[Support source]
\label{rem:three-columns-verdier}
The finite set~$S$ is the support of the family model produced by the
incidence-compatible collision calculation.  PBW and
Orlik--Solomon Koszulness provide the associated-graded quadratic
input; the OPE-residue lift and its global homotopy determine~$S$.
A perfect degree-$d$ pairing is a separate map and supplies a degree
reflection after its adjoint is a quasi-isomorphism.
De Sole--Kac variational Poisson cohomology is nonvanishing in
unboundedly many degrees for Virasoro-type Poisson vertex algebras
\cite{BDSK21}, so the concentration conclusion lives strictly on the
stated Koszul/bounded locus, with the bounded-to-chart comparison
$\chi_V^{\mathrm{bd}}$ part of $H_H(\cA;S)$.
\end{remark}

\section{The collision-depth exact triangle}
\label{sec:theorem-h-off-koszul}

Let
\[
 Q_{\mathrm H}(\cA)
 \;:=\;
 RHH_{\mathrm{ch}}(\cA)/
 F^1_{\mathrm{FM}}RHH_{\mathrm{ch}}(\cA)
\]
be the depth-zero quotient.  It is the curve-level term retained after
passing from the completed collision filtration to its depth-zero
quotient.  The long exact sequence compares its upper bound with the
positive-depth complex.  A family support datum is stronger: it
constructs the complete retract and the chart comparison in the same
ambient.

\begin{theorem}[Collision-depth tail exact sequence;
\ClaimStatusProvedHere]
\label{thm:theorem-h-off-koszul-explicit-correction}
\index{Theorem H!off Koszul locus}
Let $\cA$ be a conformal-weight-graded chiral algebra on a smooth
projective curve, with finite graded pieces, for which the completed
Fulton--MacPherson chiral Hochschild filtration is exhaustive and
Hausdorff in each conformal weight.  Fix $m\in\mathbb Z$ and assume
$Q_{\mathrm H}(\cA)\in D^{\leq m}$.  Then the exact triangle
\[
 \mathrm{KD}_{\mathrm H}^{\bullet}(\cA)
 \longrightarrow
 RHH_{\mathrm{ch}}(\cA)
 \longrightarrow
 Q_{\mathrm H}(\cA)
 \xrightarrow{+1}
\]
gives
\[
 H^m(Q_{\mathrm H}(\cA))
 \xrightarrow{\;\partial_{\mathrm H}^{m+1}(\cA)\;}
 H^{m+1}(\mathrm{KD}_{\mathrm H}^{\bullet}(\cA))
 \longrightarrow
 \ChirHoch^{m+1}(\cA)
 \longrightarrow 0
\]
and hence the first tail group
\[
 \mathfrak{o}_{\mathrm H}^{m+1}(\cA)
 :=
 \operatorname{coker}\!\left(
  \partial_{\mathrm H}^{m+1}(\cA)\colon
  H^m(Q_{\mathrm H}(\cA))
  \longrightarrow
  H^{m+1}(\mathrm{KD}_{\mathrm H}^{\bullet}(\cA))
 \right)
 \cong
 \ChirHoch^{m+1}(\cA),
\]
and, for every $n\geq m+2$,
\[
 \ChirHoch^n(\cA)
 \;\cong\;
 H^n(\mathrm{KD}_{\mathrm H}^{\bullet}(\cA)).
\]
Thus an upper support bound~$N\geq m$ follows when
$\mathfrak{o}_{\mathrm H}^{m+1}(\cA)=0$ for $N=m$ and
$H^n(\mathrm{KD}_{\mathrm H}^{\bullet}(\cA))=0$ for every $n>N$.
The specialization $m=2$ begins with the boundary
$H^2(Q_{\mathrm H})\to H^3(\mathrm{KD}_{\mathrm H})$; it is one
family calculation among the possible cutoffs.
\end{theorem}

\begin{proof}
The definition of $\mathrm{KD}_{\mathrm H}^{\bullet}$ gives the
displayed exact triangle. The long exact cohomology sequence contains
\[
H^m(Q_{\mathrm H})\to H^{m+1}(\mathrm{KD}_{\mathrm H})
\to H^{m+1}(RHH_{\mathrm{ch}})\to H^{m+1}(Q_{\mathrm H}),
\]
and the final group is zero by $Q_{\mathrm H}\in D^{\leq m}$.
For $n\geq m+2$ both $H^{n-1}(Q_{\mathrm H})$ and
$H^n(Q_{\mathrm H})$
vanish, so the map
$H^n(\mathrm{KD}_{\mathrm H})\to H^n(RHH_{\mathrm{ch}})$ is an
isomorphism. Since
$H^n(RHH_{\mathrm{ch}})=\ChirHoch^n(\cA)$, the support statements
follow.  The displayed cokernel is the first permanent collision
class remaining after the depth-zero boundary map.
\end{proof}

\begin{lemma}[Degree-three boundary for a depth-zero cutoff at two;
\ClaimStatusProvedHere]
\label{lem:theorem-h-degree-three-boundary}
\index{Theorem H!degree-three boundary obstruction}
In Theorem~\ref{thm:theorem-h-off-koszul-explicit-correction}, set
$m=2$.  Every class of $\ChirHoch^3(\cA)$ is then the image of a
positive collision-depth class modulo the boundary
$\partial_{\mathrm H}^{3}(\cA)\colon H^2(Q_{\mathrm H}(\cA))\to
H^3(\mathrm{KD}_{\mathrm H}^{\bullet}(\cA))$, and
$\mathfrak{o}_{\mathrm H}^{3}(\cA)\cong\ChirHoch^3(\cA)$.
If a family support datum has $3\in S$, this group is a retained
column of Theorem~H.  The bounded Virasoro datum is the basic example:
$S=\{0,2,3\}$ and $\dim\ChirHoch^3=1$ after the
bounded-to-chart quasi-isomorphism.
\end{lemma}

\begin{proof}
The cokernel identification is the specialization $m=2$ of
Theorem~\ref{thm:theorem-h-off-koszul-explicit-correction}.  The
Virasoro assertion is Theorem~\ref{thm:main-koszul-hoch} together
with the comparison map $\chi_{\operatorname{Vir}_c}^{\mathrm{bd}}$.
\end{proof}

\begin{corollary}[Exact Hilbert tail for the cutoff $m=2$;
\ClaimStatusProvedHere]
\label{cor:concentration-iff-defect-zero}
In the situation of
Theorem~\ref{thm:theorem-h-off-koszul-explicit-correction}, set
$m=2$.  Assume that all displayed conformal-weight graded dimensions
are finite and that the map $RHH_{\mathrm{ch}}(\cA)\to
Q_{\mathrm H}(\cA)$ induces the stated identifications in degrees
$0,1,2$ with $Z(\cA)$, $\ChirHoch^1(\cA)$, and
$Z(\cA^!)^\vee\otimes\omega_X$. Then
\[
 P_{\cA}(t,q)
 =
 \HS_{Z(\cA)}(q)
 +\HS_{\ChirHoch^1(\cA)}(q)t
 +\HS_{Z(\cA^!)}(q)t^2
 +\dim_q\mathfrak{o}_{\mathrm H}^{3}(\cA)t^3
 +\sum_{n\geq4}
  \dim_q H^n(\mathrm{KD}_{\mathrm H}^{\bullet}(\cA))\,t^n .
\]
The $t^3$ coefficient is the cokernel dimension of
$\partial_{\mathrm H}^{3}(\cA)$; it equals
$\dim_q H^3(\mathrm{KD}_{\mathrm H}^{\bullet}(\cA))$ when the
boundary map is zero.  The tail beginning in degree~$3$ remains valid
with the first three coefficients replaced by the actual
cohomology of $RHH_{\mathrm{ch}}(\cA)$.
\end{corollary}

\begin{proof}
Theorem~\ref{thm:theorem-h-off-koszul-explicit-correction} identifies
$\ChirHoch^3(\cA)$ with the cokernel
$\mathfrak{o}_{\mathrm H}^{3}(\cA)$ and identifies
$\ChirHoch^n(\cA)$ with
$H^n(\mathrm{KD}_{\mathrm H}^{\bullet}(\cA))$ for $n\geq4$. The
low-column hypothesis supplies the coefficients in degrees $0,1,2$.
\end{proof}

\begin{proposition}[Compatible finite-window retracts of the
positive-depth complex; \ClaimStatusProvedHere]
\label{prop:theorem-h-finite-window-kdh-retracts}
\index{Theorem H!finite-window KDH deformation retracts}
\index{Koszul defect complex!finite-window deformation retracts}
\textup{[Type signature: Open quadrant; completed positive-depth
chiral Hochschild chain presentation; Beilinson level~$3$; hypothesis
package: finite positive-depth windows, compatible deformation
retracts, a finite support set~$S$, and surjective weight
transitions.]}
Let $\cA$ satisfy the hypotheses of
Theorem~\ref{thm:theorem-h-off-koszul-explicit-correction}.  For the
positive-depth Hochschild defect complex set
\[
 K_N^\bullet
 :=\mathrm{KD}_{\mathrm H}^{\bullet}(\cA)/
 F_{>N}^{\mathrm{wt}}\mathrm{KD}_{\mathrm H}^{\bullet}(\cA).
\]
Suppose that completion in the conformal-weight/PBW topology gives
\[
 \mathrm{KD}_{\mathrm H}^{\bullet}(\cA)
 \;\cong\;
 \varprojlim_N K_N^\bullet
\]
and that the $K_N^\bullet$ are finite-dimensional cochain complexes
with surjective truncation maps
\(\pi_{N+1,N}\colon K_{N+1}^\bullet\twoheadrightarrow K_N^\bullet\).
Fix a finite set $S\subset\mathbb Z$.  Assume that each window
carries a finite-dimensional subcomplex \(L_N^\bullet\), whose terms
occur in degrees belonging to~$S$, chain maps
\[
 i_N\colon L_N^\bullet\longrightarrow K_N^\bullet,
 \qquad
 p_N\colon K_N^\bullet\longrightarrow L_N^\bullet,
\]
and a degree \(-1\) map
\(h_N\colon K_N^\bullet\to K_N^{\bullet-1}\) such that
\[
 p_Ni_N=\mathrm{id}_{L_N},\qquad
 d_Nh_N+h_Nd_N=\mathrm{id}_{K_N}-i_Np_N.
\]
Let
\(\rho_{N+1,N}\colon L_{N+1}^\bullet\to L_N^\bullet\)
be chain maps satisfying
\[
 \pi_{N+1,N}i_{N+1}=i_N\rho_{N+1,N},\qquad
 \rho_{N+1,N}p_{N+1}=p_N\pi_{N+1,N},\qquad
 \pi_{N+1,N}h_{N+1}=h_N\pi_{N+1,N}.
\]
Then
\[
 H^n\!\bigl(\mathrm{KD}_{\mathrm H}^{\bullet}(\cA)\bigr)=0
 \qquad(n\notin S).
\]
If $Q_{\mathrm H}(\cA)$ is also supported in~$S$, the exact triangle
gives
\[
 \operatorname{Supp}\ChirHoch^\bullet(\cA)\subseteq S.
\]

For a given family, the proposition applies once one constructs
\(K_N^\bullet\), \(L_N^\bullet\), the maps
\((i_N,p_N,h_N)\), the transitions
\((\pi_{N+1,N},\rho_{N+1,N})\), and the comparison
\(\mathrm{KD}_{\mathrm H}^{\bullet}\cong\varprojlim_N K_N^\bullet\).
These data construct the positive-depth component of a family support
model.  The chart comparison and the compatible depth-zero retract
complete the package $H_H(\cA;S)$.  De Sole--Kac variational Poisson
cohomology is nonvanishing in unboundedly many degrees for
Virasoro-type Poisson vertex algebras \cite{BDSK21}, so the resulting
concentration lives strictly on the stated Koszul/bounded locus, with
the bounded-to-chart comparison part of that package.
\end{proposition}

\begin{proof}
Fix \(N\). Let \(x\in K_N^n\) be a cocycle with \(n\notin S\).  Since
\(L_N^n=0\), one has $p_Nx=0$.  The homotopy identity gives
\[
 x=(d_Nh_N+h_Nd_N)x=d_Nh_Nx,
\]
because \(d_Nx=0\). Hence \(H^n(K_N^\bullet)=0\) for every
\(n\notin S\).

The compatibility equations define inverse-limit maps
\[
 i\colon L^\bullet:=\varprojlim_N L_N^\bullet
 \longrightarrow K^\bullet:=\varprojlim_N K_N^\bullet,
 \qquad
 p\colon K^\bullet\longrightarrow L^\bullet,
 \qquad
 h\colon K^\bullet\longrightarrow K^{\bullet-1}.
\]
They satisfy
\[
 pi=\mathrm{id}_{L},
 \qquad
 dh+hd=\mathrm{id}_{K}-ip.
\]
Since $L^n=0$ for $n\notin S$, this identity contracts every cocycle
in $K^n$ in these degrees.  The comparison
$\mathrm{KD}_{\mathrm H}^{\bullet}(\cA)\cong K^\bullet$ therefore
gives the asserted support directly.

The maps $\pi_{N+1,N}$ are surjective maps of finite-dimensional
complexes, so the inverse system is strict Mittag--Leffler. The Milnor
exact sequence records the same conclusion at the level of
cohomology: for
each \(n\),
\[
 0\to
 \varprojlim\nolimits^1_N H^{n-1}(K_N^\bullet)
 \to
 H^n\!\left(\varprojlim_N K_N^\bullet\right)
 \to
 \varprojlim_N H^n(K_N^\bullet)
 \to 0.
\]
For \(n\notin S\), the right term vanishes by the finite-window
calculation.  The left term vanishes because every finite-dimensional
cohomology tower is Mittag--Leffler: for each fixed \(N\), the
descending chain of images
\[
 \operatorname{im}\!\left(H^{n-1}(K_M^\bullet)
 \to H^{n-1}(K_N^\bullet)\right),
 \qquad M\geq N,
\]
stabilizes inside the finite-dimensional vector space
\(H^{n-1}(K_N^\bullet)\). Therefore
\(H^n(\mathrm{KD}_{\mathrm H}^{\bullet}(\cA))=0\) for
\(n\notin S\).

If $Q_{\mathrm H}(\cA)$ is supported in~$S$, the two outer terms of
the long exact sequence vanish in every degree outside~$S$; hence the
middle term $\ChirHoch^n(\cA)$ vanishes there as well.
\end{proof}

\begin{definition}[Chiral Hochschild Massey operations;
\ClaimStatusDefinitional]
\label{def:chiral-hochschild-massey-operations}
\index{Massey product!chiral Hochschild}
\index{chiral Hochschild!Massey operations}
Let
\[
 C^\bullet_{\mathrm{ch}}(\cA,\cA)
 =
 RHH_{\mathrm{ch}}(\cA)
\]
be the brace dg algebra of curve-level chiral Hochschild cochains.
For cohomology classes
\(\alpha_1,\ldots,\alpha_m\in H^\bullet C^\bullet_{\mathrm{ch}}(\cA,\cA)\),
\(m\geq3\), whose lower products admit a defining system
\(\{b_{ij}\}_{1\leq i<j\leq m}\), the \(m\)-fold chiral Hochschild
Massey operation
\[
 \langle \alpha_1,\ldots,\alpha_m\rangle_{\mathrm{ch}}
 \subset
 H^{|\alpha_1|+\cdots+|\alpha_m|-m+2}
 C^\bullet_{\mathrm{ch}}(\cA,\cA)
\]
is the usual Kadeishvili--Stasheff class of the transferred
\(A_\infty\)/brace structure on a minimal model of
\(C^\bullet_{\mathrm{ch}}(\cA,\cA)\), with the chiral cup product and
brace insertions replacing the associative Hochschild cup products.
The operation is defined modulo the standard indeterminacy from lower
Massey products. We say that \(\cA\) has \emph{Massey-vanishing on the
Theorem~H surface} if all such operations of arity \(m\geq3\) vanish
in every finite conformal-weight/PBW window and the vanishing is
compatible with the strict Mittag--Leffler completion.
\end{definition}

\begin{proposition}[Massey, formality, and comparison firewall for
Theorem~H; \ClaimStatusConditional]
\label{prop:theorem-h-massey-formality-firewall}
\index{Theorem H!Massey firewall}
\index{Fulton--MacPherson formality!Theorem H scope}
\index{mode Hochschild!filtered comparison}
\index{Lie algebra cohomology!image comparison}
Theorem~H determines cohomological support.  A free-polynomial or
strict Gerstenhaber presentation of
$H^\bullet C^\bullet_{\mathrm{ch}}(\cA,\cA)$ requires in addition the
Massey-vanishing condition of
Definition~\ref{def:chiral-hochschild-massey-operations}, equivalently
a brace-compatible formality contraction as in
Proposition~\ref{prop:e2-formality-hochschild}, has been proved on
the same finite-window completed surface.

The Fulton--MacPherson input used in Theorem~H is the rational
Orlik--Solomon/Arnold cohomology model of the type-\(A\) collision
fibres, extended to the complex coefficient field of the chiral
algebra. It supplies the fibre Koszul-complex input after
Shelton--Yuzvinsky--Priddy and the localized residue twist.  Integral
formality and the vanishing of transferred chiral Hochschild Massey
operations require their own chain contractions. Compatibility with
conformal-weight completion means: first prove the contraction in
each finite PBW/conformal-weight window, then pass through the strict
Mittag--Leffler inverse system.

For a chosen degree-two perfect pairing, the identification
\[
 H^2(Q_{\mathrm H}(\cA))
 \cong
 Z(\cA^!)^\vee\otimes\omega_X
\]
requires the Koszul/Verdier comparison for \(\cA^!\), the perfect
pairing, and a compatible depth-zero model.  Independently,
Corollary~\ref{cor:concentration-iff-defect-zero} computes the
$m=2$ tail
\(\mathfrak{o}_{\mathrm H}^{3}\oplus
 \bigoplus_{n\geq4}H^n(\mathrm{KD}_{\mathrm H}^\bullet)\).

Finally, formal-disk mode Hochschild cochains and continuous
Lie/Gel'fand--Fuchs cochains compare with \(\ChirHoch\) only by the
filtered image maps
\(\Theta_1,\Theta_{2,\ell},\Theta_{3,\ell}\) of
Theorem~\ref{thm:three-hochschild-chain-level-agreement-low-degree}
and Theorem~\ref{thm:three-hochschild-cohomological-agreement-all-degree}.
Their quasi-isomorphism ranges are the ranges supplied by the stated
PBW, spectral-kernel, and coefficient-image hypotheses. Topological,
categorical, and product-ambient Hall Hochschild theories retain their
native source categories.
\end{proposition}

\begin{proof}
The definition of Massey-vanishing is made in the brace dg algebra
$C^\bullet_{\mathrm{ch}}(\cA,\cA)$.  A formality contraction is the
additional datum that identifies its transferred $A_\infty$/brace
operations with the strict cohomology operations; this is the input of
Proposition~\ref{prop:e2-formality-hochschild}.

The FM input in the proof of Theorem~H is the type-\(A\)
Orlik--Solomon model
\(H^\bullet(\FM_m(\C);\Q)\), the Arnold relation, and the
Shelton--Yuzvinsky/Priddy Koszul-complex acyclicity after extension of
scalars. Lemma~\ref{lem:chiral-homotopy-transport} isolates the
additional residue-twisted tensor contraction on the fibre algebra.
Finite-window construction followed by strict Mittag--Leffler passage
is exactly the completion discipline in
Theorem~\ref{thm:theorem-h-on-koszul-locus}.

The low-column statement is the hypothesis explicitly isolated in
Corollary~\ref{cor:concentration-iff-defect-zero}. The comparison with
mode and continuous Lie cohomology is the typed image comparison of
Theorems~\ref{thm:three-hochschild-chain-level-agreement-low-degree}
and~\ref{thm:three-hochschild-cohomological-agreement-all-degree}.
The source-category distinctions follow from
Convention~\ref{conv:three-hochschild} and
Remark~\ref{rem:three-hochschild-discipline}.
\end{proof}

\begin{corollary}[coHochschild transport; \ClaimStatusConditional]
\label{cor:theorem-h-cohochschild-transport}
\index{coHochschild!Theorem H transport}
Let $C=\barB^{\mathrm{ch}}(\cA)$, assume the complete
bimodule--bicomodule equivalence hypotheses of
Theorem~\ref{thm:ordered-HH-coHH-cohomology}, and let $\cA$ carry
$H_H(\cA;S)$.  Then
\[
 \operatorname{Supp}\mathrm{Coext}_{C^e}^{\bullet}
 (C_\Delta,C_\Delta)\subseteq S.
\]
The same equivalence transports the collision-depth exact sequence of
Theorem~\ref{thm:theorem-h-off-koszul-explicit-correction} to the
coHochschild self-extension complex.
\end{corollary}

\begin{proof}
Theorem~\ref{thm:ordered-HH-coHH-cohomology} identifies
$\ChirHoch^\bullet(\cA)$ with
$\mathrm{Coext}_{C^e}^{\bullet}(C_\Delta,C_\Delta)$ in the complete
ordered bar ambient.  It transports the family retract and the exact
triangle, giving the stated support.
\end{proof}

\section{Logarithmic and admissible singular quotients}
\label{sec:theorem-h-logarithmic-wp}

\begin{theorem}[Triplet logarithmic boundary; \ClaimStatusConditional]
\label{thm:theorem-h-logarithmic-w(p)}
\index{Theorem H!logarithmic W(p) boundary}
\index{triplet algebra!chiral Hochschild}
Let $p\geq2$ and let $\cW(p)$ be the triplet algebra of
Definition~\ref{def:wp-algebra}. The known local facts are:
\[
 c(\cW(p))=1-\frac{6(p-1)^2}{p},\qquad
 p_{\max}(\cW(p))=2(2p-1),\qquad
 r_{\max}(\cW(p))=\infty,
\]
with $p_{\max}$ and $r_{\max}$ as in
Remark~\ref{rem:wp-pole-depth}. The algebra is $C_2$-cofinite
\textup{(}Proposition~\ref{prop:wp-c2-cofinite}\textup{)}, has
non-free strong generation
\textup{(}Proposition~\ref{prop:wp-not-free-strong}\textup{)}, and
its chiral Koszulness is Conjecture~\ref{conj:wp-koszulness}.

A Theorem~H statement for $\cW(p)$ therefore consists of a non-free
family datum $H_H(\cW(p);S_p)$.  Its collision-depth realization
incorporates the quotient relations, $E_\infty$-completion,
finite-window perfectness, and averaging comparison.  The conclusion
is
\[
 \operatorname{Supp}\ChirHoch^\bullet(\cW(p))\subseteq S_p.
\]
De Sole--Kac variational Poisson cohomology is nonvanishing in
unboundedly many degrees for Virasoro-type Poisson vertex algebras
\cite{BDSK21}, so this concentration lives strictly on the stated
Koszul/bounded locus, with the bounded-to-chart comparison part of
$H_H(\cW(p);S_p)$.
The collision-depth exact sequence computes an upper tail after a
depth-zero cutoff~$m$ has been established.  The $C_2$-cofiniteness
and maximal $W$--$W$ pole order specify finite local input; the
residue differential and its higher-page continuation determine the
cohomological support.

Let
$E_{2,W\text{-}W}^{r,s}(\cW(p))\subset E_2^{r,s}(\cW(p))$ be the
image on the $E_2$-page of the local span of collision strata whose
OPE input is a $W$-$W$ pole of order $2(2p-1)$. Let
$E_{\infty,W\text{-}W}^{r,s}(\cW(p))$ denote the subquotient of
$E_\infty^{r,s}(\cW(p))$ represented by classes in that $E_2$-subspace
which survive all higher differentials. Under the specialized
hypothesis $Q_{\mathrm H}(\cW(p))\in D^{\leq2}$, the degree-three
triplet component is
\[
 \mathfrak{o}_{\mathrm H,W}^{3}(\cW(p))
 :=
 \operatorname{im}\!\left(
  \bigoplus_{\substack{r\geq1\\ r+s=3}}
  E_{\infty,W\text{-}W}^{r,s}(\cW(p))
  \longrightarrow
  \mathfrak{o}_{\mathrm H}^{3}(\cW(p))
 \right).
\]
The differential calculation decides the map
$E_{2,W\text{-}W}^{r,s}(\cW(p))\to
E_{\infty,W\text{-}W}^{r,s}(\cW(p))$.
\end{theorem}

\begin{proof}
The displayed central charge and pole data are the defining data and
pole-depth computation in Chapter~\ref{chap:logarithmic-w-algebras}.
$C_2$-cofiniteness and non-free strong generation are exactly
Propositions~\ref{prop:wp-c2-cofinite} and
\ref{prop:wp-not-free-strong}. Theorem~H uses the PBW chiral Koszul
input and the transported Arnold contraction; $C_2$-cofiniteness is a
finite associated-variety condition, while the family contraction is
the additional chain datum. The final assertion is then the tail exact
sequence of Theorem~\ref{thm:theorem-h-off-koszul-explicit-correction}.
The pole order fixes the local $W$--$W$ summand on the first collision
divisor; the higher differentials compute its permanent image. The definition
of $\mathfrak{o}_{\mathrm H,W}^{3}$ records the permanent component
which is retained or contracted according to the chosen support
set~$S_p$.
\end{proof}

\begin{theorem}[Admissible and minimal simple-quotient firewall;
\ClaimStatusConditional]
\label{thm:theorem-h-singular-quotient-firewall}
\index{Theorem H!admissible quotient boundary}
\index{Theorem H!minimal model boundary}
An admissible simple quotient carries its own family datum
$H_H(L_k(\fg);S_{k,\fg})$.  Rationality, $C_2$-cofiniteness, and the
ordinary module category supply finite representation-theoretic input.
The curve-level chiral bar Ext complex is the additional
factorization calculation
\textup{(}Remark~\ref{rem:admissible-rationality-finiteness}\textup{)}.
For $\fg=\mathfrak{sl}_2$ chiral Koszulness at admissible level is
recorded in Remark~\ref{rem:admissible-koszul-status}; a support
conclusion follows after constructing the completion, collision
retract, chart comparison, and averaging maps in
$H_H(L_k(\mathfrak{sl}_2);S_{k,2})$. For
$\mathfrak{sl}_3$ at non-degenerate admissible level
$k=-3+p/q$ with coprime $p\geq3$, $q\geq3$, assume the finite Li--bar
package of Theorem~\ref{thm:admissible-sl3-non-koszul-qge3}, in which
the reduced bar-relevant model contains the Cartan factors
\[
 \mathbb C[H_1]/(H_1^{d_C})\otimes
 \mathbb C[H_2]/(H_2^{d_C}),
 \qquad d_C=(p-1)q,
\]
and the root factors have exponent $d_R=(p-2)q$. Assume also a
convergent completed FM Hochschild filtration, the specialized cutoff
$Q_{\mathrm H}(L_k(\mathfrak{sl}_3))\in \mathrm D^{\leq2}$, and a comparison map
from the quotient Li--bar page to the FM Hochschild collision page,
\[
 \Phi_{\mathrm H}^{\mathrm{adm}}\colon
 E_2^{2,*}\bigl(\barB(L_k(\mathfrak{sl}_3))\bigr)
 \longrightarrow
 \bigoplus_{r+s\geq3}
 E_2^{r,s}\!\left(RHH_{\mathrm{ch}}(L_k(\mathfrak{sl}_3))\right),
\]
compatible with the residue differential $d_1$. The Cartan classes
\[
 [H_1], [H_2]\in E_2^{2,*}\bigl(\barB(L_k(\mathfrak{sl}_3))\bigr),
 \qquad [\mathfrak h,\mathfrak h]=0,
\]
then define the admissible Hochschild obstruction
\[
 \mathfrak{o}_{\mathrm H,\mathrm{Cartan}}^{\ge3}(L_k(\mathfrak{sl}_3))
 :=
 \operatorname{im}\!\left(
  \Phi_{\mathrm H}^{\mathrm{adm}}\!
  \left(\operatorname{span}\{[H_1],[H_2]\}\right)
  \longrightarrow
  \mathfrak{o}_{\mathrm H,1}^{\ge3}(L_k(\mathfrak{sl}_3))
 \right),
\]
and its degree-$3$ permanent component
$\mathfrak{o}_{\mathrm H,\mathrm{Cartan}}^{3}$ is the image after
survival to $E_\infty$ followed by the edge map of
Theorem~\ref{thm:theorem-h-off-koszul-explicit-correction}. This is a
named component of $\mathfrak{o}_{\mathrm H}^{3}$.  The full group is
obtained by including every positive-depth class that survives to the
same edge map.  A support set contained in $\{0,1,2\}$ would require
\[
 \mathfrak{o}_{\mathrm H}^{3}(L_k(\mathfrak{sl}_3))=0,\qquad
 H^n(\mathrm{KD}_{\mathrm H}^{\bullet}(L_k(\mathfrak{sl}_3)))=0
 \quad(n\geq4),
\]
Together with the depth-zero calculation, a family datum may instead
retain any of these groups in its actual support set. For
$\mathrm{rk}(\fg)\geq3$ the analogous rank obstruction is
Conjecture~\ref{conj:admissible-koszul-rank-obstruction}. For
Virasoro minimal models and principal $\mathcal W$-algebras at
admissible levels, the concentration problem is the non-generic
class-\(\mathsf M\) problem of
Conjecture~\ref{conj:hochschild-concentration-class-M-non-generic}.
\end{theorem}

\begin{proof}
The cited admissible-level analysis separates ordinary representation
theory from the chiral/factorization bar complex. Theorem~H requires
the latter: completion, the family collision retract, chart
comparison, and averaging. Rationality and $C_2$-cofiniteness imply
finite ordinary data.  The support calculation additionally evaluates
the positive collision-depth complex. The rank-one and rank-two affine examples
show that the issue is mathematical rather than terminological:
rank~$1$ absorbs the admissible null vector in the quotient bar
spectral sequence, while the rank~$2$ Cartan classes survive the
$d_1$ page because the Cartan bracket vanishes. The map
$\Phi_{\mathrm H}^{\mathrm{adm}}$ is precisely the missing bridge from
bar non-Koszulness to the Hochschild collision filtration. The full
tail includes the Cartan image together with every other
positive-depth survivor. The
class-\(\mathsf M\)
statement in
Conjecture~\ref{conj:hochschild-concentration-class-M-non-generic}
records the remaining Virasoro and principal $\mathcal W$ simple
quotient problem.
\end{proof}

\section{Critical level}
\label{sec:theorem-h-critical-level}

\begin{theorem}[Critical-level centre and amplitude obstruction;
\ClaimStatusConditional]
\label{thm:theorem-h-at-critical-level}
\index{Theorem H!critical level}
\index{Feigin--Frenkel centre!Theorem H}
Let $V_{-h^\vee}(\fg)$ be the affine Kac--Moody vacuum module at
critical level for a simple Lie algebra $\fg$ of rank $r$ and
exponents $d_1,\ldots,d_r$. Its chiral centre is the Feigin--Frenkel
centre:
\[
 Z(V_{-h^\vee}(\fg))
 \;\cong\;
 \mathfrak z(\widehat{\fg})
 \;\cong\;
 \mathrm{Fun}(\mathrm{Op}_{{}^L\fg}(D)).
\]
Equivalently, as a graded commutative differential algebra it is freely
generated by the Segal--Sugawara generators $S_i$ of conformal weights
$d_i+1$ and their derivatives, so
\begin{equation}
\label{eq:ff-center-hilbert}
 \HS_{\mathfrak z(\widehat{\fg})}(q)
 =
 \prod_{i=1}^{r}\prod_{m\geq0}
 \frac{1}{1-q^{d_i+1+m}}.
\end{equation}
This computes the degree-zero chiral-centre column. A full bigraded
Hilbert series for $RHH_{\mathrm{ch}}(V_{-h^\vee}(\fg))$ is obtained
from a chain-level comparison between the critical chiral Hochschild
complex and the Beilinson--Drinfeld/Feigin--Tsygan oper complex.

The critical-first affine bar perturbation is
\[
 d_k=d_{\mathrm{crit}}+(k+h^\vee)\delta,\qquad
 d_{\mathrm{crit}}\delta+\delta d_{\mathrm{crit}}
 =
 \left[\frac{\mathsf C_{\fg}}{2h^\vee},-\right],
\]
so the obstruction to transporting the generic FM collapse across the
critical divisor is the curvature action
$[(k+h^\vee)\mathsf C_{\fg}/(2h^\vee),-]$ of
Corollary~\ref{cor:critical-level-spectral}. Equivalently, the residue
class
\[
 \mathfrak c_{\mathrm{crit}}(\fg)
 :=
 \left[\frac{\mathsf C_{\fg}}{2h^\vee},-\right]
\]
is the named obstruction to a pole-free continuation of the generic
contracting homotopy through $k=-h^\vee$. Under the separate
Beilinson--Drinfeld/Feigin--Tsygan comparison recorded in
Remark~\ref{rem:critical-level-lie-vs-chirhoch}, the continuous
Lie-cohomological factor in the critical chiral Hochschild comparison is
\[
 H^\bullet_{\mathrm{Lie,cont}}(\fg\otimes t\mathbb C[[t]];\mathbb C)
 \cong
 \Lambda(P_1,\ldots,P_r)\otimes
 \mathbb C[\Theta_1,\ldots,\Theta_r],
 \qquad
 \deg P_i=2d_i+1,\quad
 \deg\Theta_i=2(d_i+1),
\]
with the degree-zero Feigin--Frenkel centre acting as the chiral centre
when the central action is retained. Its cohomological Hilbert--Poincare
series is
\[
 \operatorname{HP}_{\mathrm{crit}}(u)
 =
 \prod_{i=1}^{r}
 \frac{1+u^{2d_i+1}}{1-u^{2(d_i+1)}}.
\]
For every $N\geq1$, the monomial $\Theta_i^N$ has cohomological degree
$2N(d_i+1)$. Thus the continuous Lie comparison has unbounded
cohomological support.  A critical-level Theorem~H statement requires
a chart comparison and a complete family model for the chart complex.
\end{theorem}

\begin{proof}
The Feigin--Frenkel theorem identifies the centre of the critical
affine vertex algebra with functions on local opers for the Langlands
dual Lie algebra; see Theorem~\ref{thm:feigin-frenkel-center} and the
formal-disc comparison in Chapter~\ref{chap:three-hochschild-unification}.
The generators have degrees $d_i+1$, and applying all derivatives
$\partial^m$ gives the product formula
\eqref{eq:ff-center-hilbert}. The final sentence is the object
firewall: a centre computation is the $0$th chiral Hochschild column,
whereas a bigraded Hochschild Hilbert series requires all cochain
degrees and all differentials. The displayed perturbation identity is
Corollary~\ref{cor:critical-level-spectral}; it is the local
calculation showing why the generic differential cannot be continued
through the critical divisor as the localized bar-concentration
collapse used in
Theorem~H. The last displayed algebra is the Lie-cohomological factor
from Remark~\ref{rem:critical-level-lie-vs-chirhoch}, with the
exponents rewritten as $d_i=m_i$; its source is the continuous Lie
comparison, while \eqref{eq:ff-center-hilbert} is the degree-zero
centre.
\end{proof}

\begin{remark}[Critical level and amplitude]
\label{rem:critical-level-separately-parametrised}
At critical level the family datum has three additional inputs: a
replacement for the Sugawara contraction, a treatment of the
curvature-residue class $\mathfrak c_{\mathrm{crit}}(\fg)$, and a
critical diagonal Ext comparison. The Hilbert series
\eqref{eq:ff-center-hilbert} has finite coefficients in each
conformal weight and determines the degree-zero column.  The remaining
columns are computed by the critical family model.
\end{remark}

\section{Boundary firewall for Theorem~H}
\label{sec:theorem-h-unified}

\begin{theorem}[Boundary firewall for Theorem~H; \ClaimStatusConditional]
\label{thm:theorem-h-platonic-unified}
\index{Theorem H!boundary firewall}
For conformal-weight-graded chiral algebras in the standard landscape
on a smooth projective curve, Theorem~H has the following typed
realizations.
\begin{enumerate}[label=\textup{(\alph*)}]
\item \emph{Completed PBW Koszul chart.}
A family datum $H_H(\cA;S)$ gives
$\operatorname{Supp}\ChirHoch^\bullet(\cA)\subseteq S$ and the
Hilbert polynomial \eqref{eq:theorem-h-on-koszul}.  The bounded
Virasoro and rank-one even superboson calculations give the benchmark
sets $\{0,2,3\}$ and $\{0,1\}$ through their bounded-to-chart maps.
De Sole--Kac variational Poisson cohomology is nonvanishing in
unboundedly many degrees for Virasoro-type Poisson vertex algebras
\cite{BDSK21}, so the concentration lives strictly on the stated
Koszul/bounded locus, with those bounded-to-chart maps part of the
hypothesis package.
Corollary~\ref{cor:theorem-h-cohochschild-transport} carries the same
set~$S$ to the diagonal coHochschild model.
\item \emph{Collision-depth calculation.}
For $Q_{\mathrm H}(\cA)\in D^{\leq m}$, the exact sequence of
Theorem~\ref{thm:theorem-h-off-koszul-explicit-correction} begins with
$\mathfrak{o}_{\mathrm H}^{m+1}(\cA)$ and identifies all degrees
$n\geq m+2$ with the cohomology of
$\mathrm{KD}_{\mathrm H}^{\bullet}(\cA)$.  Proposition~\ref{prop:theorem-h-finite-window-kdh-retracts}
constructs a prescribed finite support~$S$ from compatible windows,
transition maps, and triples $(i_N,p_N,h_N)$.
\item \emph{Logarithmic, admissible, and minimal simple quotients.}
Each quotient carries its own $H_H(\cA_{\mathrm{quot}};S_{\mathrm{quot}})$.
Rationality, $C_2$-cofiniteness, and the ordinary module category
supply finite input to that construction. The
$\mathfrak{sl}_3$, $q\geq3$ surface carries the Cartan
bar obstruction
$\mathfrak{o}_{\mathrm H,\mathrm{Cartan}}^{\ge3}$ and its
degree-three component after the comparison map
$\Phi_{\mathrm H}^{\mathrm{adm}}$ is supplied.  This component sits in
the full group
$\mathfrak{o}_{\mathrm H}^{3}$. Triplet logarithmic algebras carry the
$W$-$W$ collision component $\mathfrak{o}_{\mathrm H,W}^{3}$.
For class-\(\mathsf M\) objects, Proposition~\ref{prop:chiral-positselski-raw-direct-sum-class-M-false}
identifies the product/direct-sum separation of state families.  The
ordinary raw counit is governed by the open comparison cone of
Remark~\ref{rem:class-m-raw-comparison-open}; the completed comparison
uses its stated CP1--CP3 package.
\item \emph{Critical affine level.}
The Feigin--Frenkel theorem gives the degree-zero centre
\eqref{eq:ff-center-hilbert}. The critical perturbation is governed
by $d_k=d_{\mathrm{crit}}+(k+h^\vee)\delta$ and the residue class
$\mathfrak c_{\mathrm{crit}}(\fg)$; under the separate
BD/Feigin--Tsygan comparison, the Lie-cohomological factor has
Hilbert--Poincare series $\operatorname{HP}_{\mathrm{crit}}(u)$ with
classes in arbitrarily large cohomological degree.
\end{enumerate}
These statements have curve-level chiral Hochschild source and target.
Comparisons with $\mathrm{THH}$, categorical Hochschild cochains,
mode-algebra cochains, and product-ambient Hall cochains are explicit
functors between their respective source categories.  K3 and
higher-dimensional Calabi--Yau amplitudes enter through such a
comparison.
\end{theorem}

\begin{proof}
Item~(a) is Theorem~\ref{thm:theorem-h-on-koszul-locus} and the two
Bakalov--De Sole--Kac calculations. Item~(b) is the exact triangle
defining $\mathrm{KD}_{\mathrm H}$ and the long-exact-sequence argument of
Theorem~\ref{thm:theorem-h-off-koszul-explicit-correction}. Item~(c)
is Theorem~\ref{thm:theorem-h-singular-quotient-firewall} and
Theorem~\ref{thm:theorem-h-logarithmic-w(p)}. Item~(d) is
Theorem~\ref{thm:theorem-h-at-critical-level}. The source-category
typing is Convention~\ref{conv:three-hochschild} together with
Remark~\ref{rem:three-hochschild-discipline}.
\end{proof}
